\subsection{Comparison of valuations}

Let K be a field and A a valuation ring of K. For every subring B of K containing A, \(U(A) \subset U(B)\) . Then there is a canonical homomorphism \(\lambda\) of \(\Gamma_{A} = K^{*}/U(A)\) onto \(\Gamma_{B} = K^{*}/U(B)\) , whose kernel is \(U(B)/U(A)\) . Then, letting \(v_{A}\) and \(v_{B}\) denote the canonical valuations on K defined by A and B ( \(\S 3\) , no. 2),

\[
v _ {\mathrm{B}} = \lambda \circ v _ {\mathrm{A}}.\tag{1}
\]

As \(A \subset B\) , \(\lambda\) maps the positive elements of \(\Gamma_{A}\) to positive elements of \(\Gamma_{B}\) and hence is increasing. Therefore (Proposition 3) the kernel \(H_{B}\) of \(\lambda\) is an isolated subgroup of \(\Gamma_{A}\) and \(\lambda\) factors into \(\Gamma_{A} \to \Gamma_{A}/H_{B} \xrightarrow{\mu} \Gamma_{B}\) , where \(\mu\) is an increasing bijective homomorphism and hence an isomorphism of totally ordered groups; hence \(\Gamma_{B}\) is identified with the quotient totally ordered group \(\Gamma_{A}/H_{B}\) .

PROPOSITION 4. The mapping \(\mathbf{B} \mapsto \mathbf{H}_{\mathbf{B}}\) is an increasing bijection of the set of subrings of \(\mathbf{K}\) containing \(\mathbf{A}\) onto the set of isolated subgroups of \(\Gamma_{\mathbf{A}}\) .

Given \(H_{B}, v_{B}\) is defined up to equivalence and hence B is determined uniquely. On the other hand, let H be an isolated subgroup of \(\Gamma_{A}\) ; considering \(\Gamma_{A}/H\) as a totally ordered group (Proposition 3), the composite mapping

\[
\mathrm{K} ^ {*} \xrightarrow {v _ {\mathrm{A}}} \Gamma_ {\mathrm{A}} \longrightarrow \Gamma_ {\mathrm{A}} / \mathrm{H}
\]

is a valuation on K whose ring contains A.

Remark. Under the above hypotheses, let \(f\) denote the canonical homomorphism of \(\mathbf{B}\) onto \(\kappa(\mathbf{B})\) and \(\mathbf{A}' = f(\mathbf{A})\) ; it is a valuation ring of \(\kappa(\mathbf{B})\) (Proposition 2, no. 1). Then \(\overline{f}(\kappa(\mathbf{B})^*) = \mathbf{U}(\mathbf{B}), \overline{f}(\mathbf{A}') = \mathbf{A}, \overline{f}(\mathfrak{m}(\mathbf{A}')) = \mathfrak{m}(\mathbf{A})\) , hence

\[
\overline {{f}} (\mathbf {U} (\mathbf {A} ^ {\prime})) = \mathbf {U} (\mathbf {A}).
\]

Then there is a canonical isomorphism of \(\mathbf{U}(\mathbf{B}) / \mathbf{U}(\mathbf{A})\) onto \(\kappa (\mathbf{B})^{*} / \mathbf{U}(\mathbf{A}^{\prime}) = \Gamma_{\mathbf{A}}\) . The exact sequence

\[
0 \rightarrow \mathrm{U(B)} / \mathrm{U(A)} \rightarrow \Gamma_ {\mathrm{A}} \rightarrow \Gamma_ {\mathrm{B}} \rightarrow 0
\]

then gives an exact sequence

\[
0 \rightarrow \Gamma_ {A ^ {\prime}} \rightarrow \Gamma_ {A} \rightarrow \Gamma_ {B} \rightarrow 0.
\]

Example. Let \(k\) be a field,

\[
\mathrm{E} = k (\mathrm{X}) \quad \text { and } \quad \mathrm{K} = k (\mathrm{X}, \mathrm{Y}) = \mathrm{E} (\mathrm{Y})
\]

(X, Y indeterminates). Let \(B = E[Y]_{(Y)}\) be the valuation ring of K defined by the extremal element Y of the principal ideal domain E{[}Y{]} (§ 1, no. 4, Proposition 3). The residue field \(\kappa(\mathbf{B})\) is canonically identified with \(E[Y]/(Y) = E\) . Similarly, let \(A' = k[X]_{(X)}\) be the valuation ring of \(E = k(X)\) defined by the extremal element X of \(k[X]\) . Denoting by \(h_{B}\) the place of E associated with B and writing \(A = \overline{h_{B}}(A')\) , a valuation ring A of K is defined which is contained in B and \(\kappa(A) = \kappa(A') = k\) . The canonical place \(h_{A}: K \to k\) can be described as follows: if \(f(X, Y)\) is an element of K, then we first put Y = 0 in f (which gives an element of \(\hat{E} = k(X)^{-}\) ), then X = 0 in the result obtained. The groups \(\Gamma_{A'}\) and \(\Gamma_{B}\) are canonically isomorphic to Z (§ 3, no. 4, Example 4). *It is not difficult to show (cf.~§ 10, no. 2, Lemma 2) that the group \(\Gamma_{A}\) is isomorphic to the lexicographic product \(Z \times Z\) and that the valuation \(v_{A}\) is equivalent to the valuation defined in § 3, no. 4, end of Example 6.*

